\documentclass{article} % For LaTeX2e
\usepackage[T1]{fontenc}
\usepackage{iclr2025_conference,times}
\iclrfinalcopy % not anonymous: an internal report, not a double-blind submission

% Optional math commands from https://github.com/goodfeli/dlbook_notation.
%%%%% NEW MATH DEFINITIONS %%%%%

\usepackage{amsmath,amsfonts,bm}

% Mark sections of captions for referring to divisions of figures
\newcommand{\figleft}{{\em (Left)}}
\newcommand{\figcenter}{{\em (Center)}}
\newcommand{\figright}{{\em (Right)}}
\newcommand{\figtop}{{\em (Top)}}
\newcommand{\figbottom}{{\em (Bottom)}}
\newcommand{\captiona}{{\em (a)}}
\newcommand{\captionb}{{\em (b)}}
\newcommand{\captionc}{{\em (c)}}
\newcommand{\captiond}{{\em (d)}}

% Highlight a newly defined term
\newcommand{\newterm}[1]{{\bf #1}}


% Figure reference, lower-case.
\def\figref#1{figure~\ref{#1}}
% Figure reference, capital. For start of sentence
\def\Figref#1{Figure~\ref{#1}}
\def\twofigref#1#2{figures \ref{#1} and \ref{#2}}
\def\quadfigref#1#2#3#4{figures \ref{#1}, \ref{#2}, \ref{#3} and \ref{#4}}
% Section reference, lower-case.
\def\secref#1{section~\ref{#1}}
% Section reference, capital.
\def\Secref#1{Section~\ref{#1}}
% Reference to two sections.
\def\twosecrefs#1#2{sections \ref{#1} and \ref{#2}}
% Reference to three sections.
\def\secrefs#1#2#3{sections \ref{#1}, \ref{#2} and \ref{#3}}
% Reference to an equation, lower-case.
\def\eqref#1{equation~\ref{#1}}
% Reference to an equation, upper case
\def\Eqref#1{Equation~\ref{#1}}
% A raw reference to an equation---avoid using if possible
\def\plaineqref#1{\ref{#1}}
% Reference to a chapter, lower-case.
\def\chapref#1{chapter~\ref{#1}}
% Reference to an equation, upper case.
\def\Chapref#1{Chapter~\ref{#1}}
% Reference to a range of chapters
\def\rangechapref#1#2{chapters\ref{#1}--\ref{#2}}
% Reference to an algorithm, lower-case.
\def\algref#1{algorithm~\ref{#1}}
% Reference to an algorithm, upper case.
\def\Algref#1{Algorithm~\ref{#1}}
\def\twoalgref#1#2{algorithms \ref{#1} and \ref{#2}}
\def\Twoalgref#1#2{Algorithms \ref{#1} and \ref{#2}}
% Reference to a part, lower case
\def\partref#1{part~\ref{#1}}
% Reference to a part, upper case
\def\Partref#1{Part~\ref{#1}}
\def\twopartref#1#2{parts \ref{#1} and \ref{#2}}

\def\ceil#1{\lceil #1 \rceil}
\def\floor#1{\lfloor #1 \rfloor}
\def\1{\bm{1}}
\newcommand{\train}{\mathcal{D}}
\newcommand{\valid}{\mathcal{D_{\mathrm{valid}}}}
\newcommand{\test}{\mathcal{D_{\mathrm{test}}}}

\def\eps{{\epsilon}}


% Random variables
\def\reta{{\textnormal{$\eta$}}}
\def\ra{{\textnormal{a}}}
\def\rb{{\textnormal{b}}}
\def\rc{{\textnormal{c}}}
\def\rd{{\textnormal{d}}}
\def\re{{\textnormal{e}}}
\def\rf{{\textnormal{f}}}
\def\rg{{\textnormal{g}}}
\def\rh{{\textnormal{h}}}
\def\ri{{\textnormal{i}}}
\def\rj{{\textnormal{j}}}
\def\rk{{\textnormal{k}}}
\def\rl{{\textnormal{l}}}
% rm is already a command, just don't name any random variables m
\def\rn{{\textnormal{n}}}
\def\ro{{\textnormal{o}}}
\def\rp{{\textnormal{p}}}
\def\rq{{\textnormal{q}}}
\def\rr{{\textnormal{r}}}
\def\rs{{\textnormal{s}}}
\def\rt{{\textnormal{t}}}
\def\ru{{\textnormal{u}}}
\def\rv{{\textnormal{v}}}
\def\rw{{\textnormal{w}}}
\def\rx{{\textnormal{x}}}
\def\ry{{\textnormal{y}}}
\def\rz{{\textnormal{z}}}

% Random vectors
\def\rvepsilon{{\mathbf{\epsilon}}}
\def\rvtheta{{\mathbf{\theta}}}
\def\rva{{\mathbf{a}}}
\def\rvb{{\mathbf{b}}}
\def\rvc{{\mathbf{c}}}
\def\rvd{{\mathbf{d}}}
\def\rve{{\mathbf{e}}}
\def\rvf{{\mathbf{f}}}
\def\rvg{{\mathbf{g}}}
\def\rvh{{\mathbf{h}}}
\def\rvu{{\mathbf{i}}}
\def\rvj{{\mathbf{j}}}
\def\rvk{{\mathbf{k}}}
\def\rvl{{\mathbf{l}}}
\def\rvm{{\mathbf{m}}}
\def\rvn{{\mathbf{n}}}
\def\rvo{{\mathbf{o}}}
\def\rvp{{\mathbf{p}}}
\def\rvq{{\mathbf{q}}}
\def\rvr{{\mathbf{r}}}
\def\rvs{{\mathbf{s}}}
\def\rvt{{\mathbf{t}}}
\def\rvu{{\mathbf{u}}}
\def\rvv{{\mathbf{v}}}
\def\rvw{{\mathbf{w}}}
\def\rvx{{\mathbf{x}}}
\def\rvy{{\mathbf{y}}}
\def\rvz{{\mathbf{z}}}

% Elements of random vectors
\def\erva{{\textnormal{a}}}
\def\ervb{{\textnormal{b}}}
\def\ervc{{\textnormal{c}}}
\def\ervd{{\textnormal{d}}}
\def\erve{{\textnormal{e}}}
\def\ervf{{\textnormal{f}}}
\def\ervg{{\textnormal{g}}}
\def\ervh{{\textnormal{h}}}
\def\ervi{{\textnormal{i}}}
\def\ervj{{\textnormal{j}}}
\def\ervk{{\textnormal{k}}}
\def\ervl{{\textnormal{l}}}
\def\ervm{{\textnormal{m}}}
\def\ervn{{\textnormal{n}}}
\def\ervo{{\textnormal{o}}}
\def\ervp{{\textnormal{p}}}
\def\ervq{{\textnormal{q}}}
\def\ervr{{\textnormal{r}}}
\def\ervs{{\textnormal{s}}}
\def\ervt{{\textnormal{t}}}
\def\ervu{{\textnormal{u}}}
\def\ervv{{\textnormal{v}}}
\def\ervw{{\textnormal{w}}}
\def\ervx{{\textnormal{x}}}
\def\ervy{{\textnormal{y}}}
\def\ervz{{\textnormal{z}}}

% Random matrices
\def\rmA{{\mathbf{A}}}
\def\rmB{{\mathbf{B}}}
\def\rmC{{\mathbf{C}}}
\def\rmD{{\mathbf{D}}}
\def\rmE{{\mathbf{E}}}
\def\rmF{{\mathbf{F}}}
\def\rmG{{\mathbf{G}}}
\def\rmH{{\mathbf{H}}}
\def\rmI{{\mathbf{I}}}
\def\rmJ{{\mathbf{J}}}
\def\rmK{{\mathbf{K}}}
\def\rmL{{\mathbf{L}}}
\def\rmM{{\mathbf{M}}}
\def\rmN{{\mathbf{N}}}
\def\rmO{{\mathbf{O}}}
\def\rmP{{\mathbf{P}}}
\def\rmQ{{\mathbf{Q}}}
\def\rmR{{\mathbf{R}}}
\def\rmS{{\mathbf{S}}}
\def\rmT{{\mathbf{T}}}
\def\rmU{{\mathbf{U}}}
\def\rmV{{\mathbf{V}}}
\def\rmW{{\mathbf{W}}}
\def\rmX{{\mathbf{X}}}
\def\rmY{{\mathbf{Y}}}
\def\rmZ{{\mathbf{Z}}}

% Elements of random matrices
\def\ermA{{\textnormal{A}}}
\def\ermB{{\textnormal{B}}}
\def\ermC{{\textnormal{C}}}
\def\ermD{{\textnormal{D}}}
\def\ermE{{\textnormal{E}}}
\def\ermF{{\textnormal{F}}}
\def\ermG{{\textnormal{G}}}
\def\ermH{{\textnormal{H}}}
\def\ermI{{\textnormal{I}}}
\def\ermJ{{\textnormal{J}}}
\def\ermK{{\textnormal{K}}}
\def\ermL{{\textnormal{L}}}
\def\ermM{{\textnormal{M}}}
\def\ermN{{\textnormal{N}}}
\def\ermO{{\textnormal{O}}}
\def\ermP{{\textnormal{P}}}
\def\ermQ{{\textnormal{Q}}}
\def\ermR{{\textnormal{R}}}
\def\ermS{{\textnormal{S}}}
\def\ermT{{\textnormal{T}}}
\def\ermU{{\textnormal{U}}}
\def\ermV{{\textnormal{V}}}
\def\ermW{{\textnormal{W}}}
\def\ermX{{\textnormal{X}}}
\def\ermY{{\textnormal{Y}}}
\def\ermZ{{\textnormal{Z}}}

% Vectors
\def\vzero{{\bm{0}}}
\def\vone{{\bm{1}}}
\def\vmu{{\bm{\mu}}}
\def\vtheta{{\bm{\theta}}}
\def\va{{\bm{a}}}
\def\vb{{\bm{b}}}
\def\vc{{\bm{c}}}
\def\vd{{\bm{d}}}
\def\ve{{\bm{e}}}
\def\vf{{\bm{f}}}
\def\vg{{\bm{g}}}
\def\vh{{\bm{h}}}
\def\vi{{\bm{i}}}
\def\vj{{\bm{j}}}
\def\vk{{\bm{k}}}
\def\vl{{\bm{l}}}
\def\vm{{\bm{m}}}
\def\vn{{\bm{n}}}
\def\vo{{\bm{o}}}
\def\vp{{\bm{p}}}
\def\vq{{\bm{q}}}
\def\vr{{\bm{r}}}
\def\vs{{\bm{s}}}
\def\vt{{\bm{t}}}
\def\vu{{\bm{u}}}
\def\vv{{\bm{v}}}
\def\vw{{\bm{w}}}
\def\vx{{\bm{x}}}
\def\vy{{\bm{y}}}
\def\vz{{\bm{z}}}

% Elements of vectors
\def\evalpha{{\alpha}}
\def\evbeta{{\beta}}
\def\evepsilon{{\epsilon}}
\def\evlambda{{\lambda}}
\def\evomega{{\omega}}
\def\evmu{{\mu}}
\def\evpsi{{\psi}}
\def\evsigma{{\sigma}}
\def\evtheta{{\theta}}
\def\eva{{a}}
\def\evb{{b}}
\def\evc{{c}}
\def\evd{{d}}
\def\eve{{e}}
\def\evf{{f}}
\def\evg{{g}}
\def\evh{{h}}
\def\evi{{i}}
\def\evj{{j}}
\def\evk{{k}}
\def\evl{{l}}
\def\evm{{m}}
\def\evn{{n}}
\def\evo{{o}}
\def\evp{{p}}
\def\evq{{q}}
\def\evr{{r}}
\def\evs{{s}}
\def\evt{{t}}
\def\evu{{u}}
\def\evv{{v}}
\def\evw{{w}}
\def\evx{{x}}
\def\evy{{y}}
\def\evz{{z}}

% Matrix
\def\mA{{\bm{A}}}
\def\mB{{\bm{B}}}
\def\mC{{\bm{C}}}
\def\mD{{\bm{D}}}
\def\mE{{\bm{E}}}
\def\mF{{\bm{F}}}
\def\mG{{\bm{G}}}
\def\mH{{\bm{H}}}
\def\mI{{\bm{I}}}
\def\mJ{{\bm{J}}}
\def\mK{{\bm{K}}}
\def\mL{{\bm{L}}}
\def\mM{{\bm{M}}}
\def\mN{{\bm{N}}}
\def\mO{{\bm{O}}}
\def\mP{{\bm{P}}}
\def\mQ{{\bm{Q}}}
\def\mR{{\bm{R}}}
\def\mS{{\bm{S}}}
\def\mT{{\bm{T}}}
\def\mU{{\bm{U}}}
\def\mV{{\bm{V}}}
\def\mW{{\bm{W}}}
\def\mX{{\bm{X}}}
\def\mY{{\bm{Y}}}
\def\mZ{{\bm{Z}}}
\def\mBeta{{\bm{\beta}}}
\def\mPhi{{\bm{\Phi}}}
\def\mLambda{{\bm{\Lambda}}}
\def\mSigma{{\bm{\Sigma}}}

% Tensor
\DeclareMathAlphabet{\mathsfit}{\encodingdefault}{\sfdefault}{m}{sl}
\SetMathAlphabet{\mathsfit}{bold}{\encodingdefault}{\sfdefault}{bx}{n}
\newcommand{\tens}[1]{\bm{\mathsfit{#1}}}
\def\tA{{\tens{A}}}
\def\tB{{\tens{B}}}
\def\tC{{\tens{C}}}
\def\tD{{\tens{D}}}
\def\tE{{\tens{E}}}
\def\tF{{\tens{F}}}
\def\tG{{\tens{G}}}
\def\tH{{\tens{H}}}
\def\tI{{\tens{I}}}
\def\tJ{{\tens{J}}}
\def\tK{{\tens{K}}}
\def\tL{{\tens{L}}}
\def\tM{{\tens{M}}}
\def\tN{{\tens{N}}}
\def\tO{{\tens{O}}}
\def\tP{{\tens{P}}}
\def\tQ{{\tens{Q}}}
\def\tR{{\tens{R}}}
\def\tS{{\tens{S}}}
\def\tT{{\tens{T}}}
\def\tU{{\tens{U}}}
\def\tV{{\tens{V}}}
\def\tW{{\tens{W}}}
\def\tX{{\tens{X}}}
\def\tY{{\tens{Y}}}
\def\tZ{{\tens{Z}}}


% Graph
\def\gA{{\mathcal{A}}}
\def\gB{{\mathcal{B}}}
\def\gC{{\mathcal{C}}}
\def\gD{{\mathcal{D}}}
\def\gE{{\mathcal{E}}}
\def\gF{{\mathcal{F}}}
\def\gG{{\mathcal{G}}}
\def\gH{{\mathcal{H}}}
\def\gI{{\mathcal{I}}}
\def\gJ{{\mathcal{J}}}
\def\gK{{\mathcal{K}}}
\def\gL{{\mathcal{L}}}
\def\gM{{\mathcal{M}}}
\def\gN{{\mathcal{N}}}
\def\gO{{\mathcal{O}}}
\def\gP{{\mathcal{P}}}
\def\gQ{{\mathcal{Q}}}
\def\gR{{\mathcal{R}}}
\def\gS{{\mathcal{S}}}
\def\gT{{\mathcal{T}}}
\def\gU{{\mathcal{U}}}
\def\gV{{\mathcal{V}}}
\def\gW{{\mathcal{W}}}
\def\gX{{\mathcal{X}}}
\def\gY{{\mathcal{Y}}}
\def\gZ{{\mathcal{Z}}}

% Sets
\def\sA{{\mathbb{A}}}
\def\sB{{\mathbb{B}}}
\def\sC{{\mathbb{C}}}
\def\sD{{\mathbb{D}}}
% Don't use a set called E, because this would be the same as our symbol
% for expectation.
\def\sF{{\mathbb{F}}}
\def\sG{{\mathbb{G}}}
\def\sH{{\mathbb{H}}}
\def\sI{{\mathbb{I}}}
\def\sJ{{\mathbb{J}}}
\def\sK{{\mathbb{K}}}
\def\sL{{\mathbb{L}}}
\def\sM{{\mathbb{M}}}
\def\sN{{\mathbb{N}}}
\def\sO{{\mathbb{O}}}
\def\sP{{\mathbb{P}}}
\def\sQ{{\mathbb{Q}}}
\def\sR{{\mathbb{R}}}
\def\sS{{\mathbb{S}}}
\def\sT{{\mathbb{T}}}
\def\sU{{\mathbb{U}}}
\def\sV{{\mathbb{V}}}
\def\sW{{\mathbb{W}}}
\def\sX{{\mathbb{X}}}
\def\sY{{\mathbb{Y}}}
\def\sZ{{\mathbb{Z}}}

% Entries of a matrix
\def\emLambda{{\Lambda}}
\def\emA{{A}}
\def\emB{{B}}
\def\emC{{C}}
\def\emD{{D}}
\def\emE{{E}}
\def\emF{{F}}
\def\emG{{G}}
\def\emH{{H}}
\def\emI{{I}}
\def\emJ{{J}}
\def\emK{{K}}
\def\emL{{L}}
\def\emM{{M}}
\def\emN{{N}}
\def\emO{{O}}
\def\emP{{P}}
\def\emQ{{Q}}
\def\emR{{R}}
\def\emS{{S}}
\def\emT{{T}}
\def\emU{{U}}
\def\emV{{V}}
\def\emW{{W}}
\def\emX{{X}}
\def\emY{{Y}}
\def\emZ{{Z}}
\def\emSigma{{\Sigma}}

% entries of a tensor
% Same font as tensor, without \bm wrapper
\newcommand{\etens}[1]{\mathsfit{#1}}
\def\etLambda{{\etens{\Lambda}}}
\def\etA{{\etens{A}}}
\def\etB{{\etens{B}}}
\def\etC{{\etens{C}}}
\def\etD{{\etens{D}}}
\def\etE{{\etens{E}}}
\def\etF{{\etens{F}}}
\def\etG{{\etens{G}}}
\def\etH{{\etens{H}}}
\def\etI{{\etens{I}}}
\def\etJ{{\etens{J}}}
\def\etK{{\etens{K}}}
\def\etL{{\etens{L}}}
\def\etM{{\etens{M}}}
\def\etN{{\etens{N}}}
\def\etO{{\etens{O}}}
\def\etP{{\etens{P}}}
\def\etQ{{\etens{Q}}}
\def\etR{{\etens{R}}}
\def\etS{{\etens{S}}}
\def\etT{{\etens{T}}}
\def\etU{{\etens{U}}}
\def\etV{{\etens{V}}}
\def\etW{{\etens{W}}}
\def\etX{{\etens{X}}}
\def\etY{{\etens{Y}}}
\def\etZ{{\etens{Z}}}

% The true underlying data generating distribution
\newcommand{\pdata}{p_{\rm{data}}}
% The empirical distribution defined by the training set
\newcommand{\ptrain}{\hat{p}_{\rm{data}}}
\newcommand{\Ptrain}{\hat{P}_{\rm{data}}}
% The model distribution
\newcommand{\pmodel}{p_{\rm{model}}}
\newcommand{\Pmodel}{P_{\rm{model}}}
\newcommand{\ptildemodel}{\tilde{p}_{\rm{model}}}
% Stochastic autoencoder distributions
\newcommand{\pencode}{p_{\rm{encoder}}}
\newcommand{\pdecode}{p_{\rm{decoder}}}
\newcommand{\precons}{p_{\rm{reconstruct}}}

\newcommand{\laplace}{\mathrm{Laplace}} % Laplace distribution

\newcommand{\E}{\mathbb{E}}
\newcommand{\Ls}{\mathcal{L}}
\newcommand{\R}{\mathbb{R}}
\newcommand{\emp}{\tilde{p}}
\newcommand{\lr}{\alpha}
\newcommand{\reg}{\lambda}
\newcommand{\rect}{\mathrm{rectifier}}
\newcommand{\softmax}{\mathrm{softmax}}
\newcommand{\sigmoid}{\sigma}
\newcommand{\softplus}{\zeta}
\newcommand{\KL}{D_{\mathrm{KL}}}
\newcommand{\Var}{\mathrm{Var}}
\newcommand{\standarderror}{\mathrm{SE}}
\newcommand{\Cov}{\mathrm{Cov}}
% Wolfram Mathworld says $L^2$ is for function spaces and $\ell^2$ is for vectors
% But then they seem to use $L^2$ for vectors throughout the site, and so does
% wikipedia.
\newcommand{\normlzero}{L^0}
\newcommand{\normlone}{L^1}
\newcommand{\normltwo}{L^2}
\newcommand{\normlp}{L^p}
\newcommand{\normmax}{L^\infty}

\newcommand{\parents}{Pa} % See usage in notation.tex. Chosen to match Daphne's book.

\DeclareMathOperator*{\argmax}{arg\,max}
\DeclareMathOperator*{\argmin}{arg\,min}

\DeclareMathOperator{\sign}{sign}
\DeclareMathOperator{\Tr}{Tr}
\let\ab\allowbreak


\usepackage{hyperref}
\usepackage{url}
\usepackage{booktabs}
\usepackage{graphicx}

% Human results not yet in print as a boxed "pending" marker; see human_results.tex
% Values of the human results (calibration, audit, triage, adjudication, bias sheets).
% \HR{<id>} prints the value when it is defined here, and a boxed "pending" marker until then.
% The ids are listed in reports/paper_v4/materials/pending_human_results.md.
% Do not edit by hand: analysis/paper_v4/fill_human_results.py writes the definitions below.
\makeatletter
\newcommand{\HR}[1]{\@ifundefined{hr@#1}{\fbox{\scriptsize\textsf{pending:\,\detokenize{#1}}}}{\@nameuse{hr@#1}}}
\makeatother
% ---- values (none yet, 2026-11-20) ----


%%%%%%%%% TITLE - PLEASE UPDATE TITLE HERE
\title{My Title}

\author{Jimmy Lin \\
Rare Genomics Institute
\And
Monica S. Lam \\
Stanford University, Open Virtual Assistant Lab
}

\newcommand{\fix}{\marginpar{FIX}}
\newcommand{\new}{\marginpar{NEW}}

\begin{document}

\maketitle

\begin{abstract}
%%%%%%%%% ABSTRACT %%%%%%%%%
Abstract here.
%%%%%%%%% ABSTRACT %%%%%%%%%
\end{abstract}

\section{Introduction}
\label{sec:intro}
%%%%%%%%% INTRODUCTION %%%%%%%%%
Rare diseases together affect an estimated 3.5--5.9\% of the population \citep{Wakap2019EstimatingCumulative}, and in Europe the average time to diagnosis is close to five years \citep{Faye2024TimeDiagnosis}. Reanalysis of exome data solves further cases as new gene--disease evidence is published \citep{Wenger2017SystematicReanalysis}, with a median new-diagnosis rate of 15\% across the studies one review collected \citep{Tan2020EvaluatingSystematic}, so a diagnosis often depends on whether a laboratory can find evidence that already exists. That evidence is fragmented. Variant interpretations sit in ClinVar \citep{Landrum2017ClinvarImproving,landrum2025_clinvar} and OMIM \citep{amberger2019_omim,hamosh2021_omim}, clinical relevance of genes and variants is curated by ClinGen \citep{Rehm2015ClingenClinical}, phenotype vocabulary lives in the Human Phenotype Ontology (HPO) \citep{Robinson2008HumanPhenotype,gargano2024_hpo} that underpins deep phenotyping \citep{Robinson2012DeepPhenotyping}, disease terminology in Mondo \citep{vasilevsky2025_mondo} and Orphanet \citep{rath2012_orphanet,lucano2026_orphanet}, and a small share of individual cases in Monarch's phenopacket-store \citep{putman2024_monarch,danis2025_phenopacket_store}. Most individual patients, however, are described only in published case reports, and a laboratory querying ClinVar cannot see a causative variant that was published but never submitted. In one curated gene, 24\% of the pathogenic variants found in PubMed articles were absent from every database examined \citep{Lyons2023RareDisease}; a full-text retrieval study reports more than 60,000 likely disease-causing HGMD variants absent from ClinVar \citep{birgmeier2020_avada}; and registries that pair clinical variants with phenotypes report variants new to ClinVar \citep{Chopra2024ClinicalVariants}. Non-standard variant names in papers hide evidence from databases \citep{freeman2024_vv_literature}, every one of 52 audited manuscript submissions to one genetics journal contained gene or variant nomenclature errors \citep{lansdon2026_nomenclature_errors}, and phenotype detail in first reports of ultra-rare conditions is often superficial or absent \citep{almail2024_reporting_gaps}. Gene-discovery platforms such as the Matchmaker Exchange \citep{philippakis2015_mme} can only match what has been recorded.

Alexander disease shows the size of the gap. A 2024 systematic review and meta-analysis reports 550 published patients and 182 causative \emph{GFAP} variants \citep{grossi2024_axd_meta}; ClinVar's release of 2026-10-12 holds 76 pathogenic or likely pathogenic (P/LP) single-gene \emph{GFAP} records, and OMIM's \emph{GFAP} gene entry (frozen 2026-10-28) lists 13 allelic variants. None of the twelve genes studied here has a folder in phenopacket-store (735 gene folders and 10,809 phenopackets at commit \texttt{ca133e6}, 2026-11-04). Uncurated diseases come up faster than hand curation can follow, so what is needed is a pipeline that can be pointed at any disease and returns records whose every value a clinician can check against the paper.

Large language models (LLMs) are an obvious candidate. Surveys of LLMs in biomedicine, evidence synthesis and medical genetics report rapid progress alongside fabricated content, limited generalization and a lack of shared benchmarks \citep{Tian2023OpportunitiesChallenges,Garcia2024ReviewScientific,Changalidis2026SystematicReview}, and LLM data extraction for evidence synthesis is now the subject of its own systematic review \citep{Shankar2026PerformanceLarge}. Much LLM evaluation in rare disease has targeted diagnosis \citep{Chen2024RarebenchCan}: on about 5,200 phenopacket-derived cases from the literature, LLMs have not reached the top-1 accuracy of a traditional decision-support tool \citep{reese2026_llm_vs_exomiser}, and a meta-analysis rated all 19 system--dataset pairs at high risk of bias, mainly from data leakage \citep{nguyen2026_llm_rd_metaanalysis}. Building the case-level records themselves, with an error rate measured by people, has received less attention.

GP-RareLink asks whether SLIDERS \citep{joshi2026_sliders}, an LLM-based systematic-review method that builds question-specific evidence tables with verbatim excerpts as provenance and merges and de-conflicts extractions with a code-writing reconciliation agent, can, run with Claude, build consolidated variant--phenotype--disease records from published case reports at an accuracy we can demonstrate. We run a home-built version of the method with Claude Code helper agents, not the public SLIDERS code. For each open-access case report, Claude reads the full text and writes per-person records in which every value carries a verbatim quote and the paper's PMID; phenotypes are three-state (observed, excluded, not reported), and ``not reported'' is never collapsed into ``excluded''. A second, independent Claude read, the \emph{gatekeeper}, checks every value against its own quote and removes values the quote does not support. Variants are normalized to HGVS \citep{dendunnen2016_hgvs,hart2024_hgvs2024} on the MANE Select transcript \citep{morales2022_mane} with VariantValidator \citep{freeman2018_variantvalidator}, keeping the paper's original string; features are bound to HPO and diseases to Mondo; and the output is GA4GH Phenopackets v2 \citep{jacobsen2022_phenopacket_schema,ladewig2023_phenopackets_primer}, validated against the schema \citep{danis2023_phenopacket_tools}. The registered corpus uses open-access full text only, and open access covers only part of the literature \citep{Piwowar2018StateOa}: 421 of its 1,496 papers (28\%) had it. ClinVar and OMIM are never read by the extraction pipeline, a rule enforced by tests, so that afterwards they can score it: concordance (O2) and variants absent from both databases, called gap candidates (O3). The analysis plan, prompts and schema were frozen on 2026-10-20 before the main runs, and every later change is a dated deviation. People decide what counts: an audit of extracted values against their quotes (O1, the primary outcome, with a target of fewer than 2\% unsupported on a 250-claim audit), a triage of ClinVar misses, and adjudication of gap candidates; nothing goes to a database without the adjudication panel's approval and the PI's sign-off.

The registered study covers six diseases that span four inheritance modes and six organ systems: A20 haploinsufficiency (HA20; \emph{TNFAIP3}) \citep{zhou2016_ha20_original,aeschlimann2018_ha20_cohort,elhani2024_ha20_177,genereviews_ha20,li2026_ha20_howitreat,he2026_ha20_multicenter}, SSADH deficiency (\emph{ALDH5A1}) \citep{pearl2003_ssadh_series,didiasova2020_ssadh_update,tokatlylatzer2023_ssadh_genotype,tokatlylatzer2024_ssadh_guidelines,genereviews_ssadh}, Tangier disease (\emph{ABCA1}) \citep{puntoni2012_tangier,schaefer2016_hdl_deficiency,mercan2018_tangier_neuropathy,genereviews_tangier,koseki2021_tangier_review}, leukocyte adhesion deficiency type I (LAD-I; \emph{ITGB2}) \citep{qasim2009_lad_hsct,almarzanovoa2018_lad1_review,kambli2020_lad_india,booth2025_lad_gene_therapy}, Alexander disease (AxD; \emph{GFAP}) \citep{prust2011_axd_onset,genereviews_axd,grossi2024_axd_meta,messing2025_axd_14variants}, and Danon disease (\emph{LAMP2}) \citep{sugie2018_danon_japan,brambatti2019_danon_gender,cenacchi2020_danon_review,genereviews_danon,sugie2024_danon_history}. Natural-history comparator cohorts for the bias arm \citep{mura2021_axd_italian,gagliardi2024_axd_ukb,dibacco2019_ssadh_natural_history,boucek2011_danon_natural_history,khedro2025_danon_prospective} are held out of the extraction corpus. HA20 was the development disease on which the extraction rules were tuned, so headline pooled figures use the other five. The same frozen pipeline was also run, under a separately pre-registered extension that is a recorded deviation, on five reserve diseases: Blau syndrome (\emph{NOD2}) \citep{wouters2014_blau,rose2015_blau_cohort,wang2025_blau_review}, familial LCAT deficiency \citep{vitali2022_lcat}, GAMT deficiency \citep{stockleripsiroglu2014_gamt,biagiotti2025_gamt_review}, arterial tortuosity syndrome (\emph{SLC2A10}) \citep{beyens2018_ats}, and PAPA syndrome (\emph{PSTPIP1}) \citep{genovese2020_papa_spectrum,hajialigol2022_papa,satoh2024_pg_genetics}; and on Majeed syndrome (\emph{LPIN2}) \citep{ferguson2021_majeed} as an exploratory twelfth disease. The extension and Majeed are not part of the registered study; their results are reported beside the registered six and never pooled into them.

\paragraph{Status of this report.} This is version 4 of an internal interim report; it supersedes version 2 (2026-10-28) and version 3 (2026-11-02). It is written before any person has rated output of the frozen pipeline. Every human result is printed as a boxed placeholder that a script fills when the ratings arrive; the primary outcome, for example, will read: \HR{o1_unsupported_n} of 250 audited claims unsupported (\HR{o1_unsupported_pct}; exact 95\% CI \HR{o1_unsupported_ci}). The only human ratings of extracted claims so far come from two pilot audits on HA20 before the freeze, each on an earlier pipeline version and by one rater, and both were above the 2\% target: 6.0\% (3 of 50) and 6.3\% (3 of 48). The share of claims the gatekeeper removes is a second model reading, not an error rate.

\paragraph{Contributions.} As this report can claim them:
\begin{itemize}
\item An end-to-end, provenance-first pipeline from case reports to validated phenopackets, run with the frozen pipeline on twelve rare diseases (the registered six, five extension diseases and one exploratory disease; 665 paper-runs over 660 distinct papers). In the registered six, every one of 241,730 asserted values has a quote that a deterministic check found verbatim in its paper, and all 2,836 phenopackets are schema-valid.
\item A database-held-out evaluation design: automatic concordance with ClinVar and OMIM, and a list of gap candidates with evidence dossiers ready for expert adjudication (215 in the registered six, 214 of them live; 269 across all twelve diseases left for the panel).
\item An honest account of what is not yet known: the primary outcome waits on human raters, the pilot audits were above the 2\% target, and HGVS recall against ClinVar is low (10--35\% at tier~1 in the registered six) with its misses concentrated in a few papers.
\end{itemize}
%%%%%%%%% INTRODUCTION %%%%%%%%%

\section{Related Work}
\label{sec:related_work}
%%%%%%%%% RELATED WORK %%%%%%%%%
\subsection{Variant and phenotype extraction from the biomedical literature}

\paragraph{Variants.} tmVar~3.0 \citep{wei2022_tmvar3} is the standard recognizer and normalizer of variant mentions, and LitVar~2.0 \citep{allot2023_litvar2} indexes about 15 million variants against the papers that mention them; because LitVar links to ClinVar, those links are kept out of our extraction step. AVADA \citep{birgmeier2020_avada} retrieved pathogenic-variant evidence directly from full text and recovered 59\% of HGMD's disease-causing variants. LLM pipelines now work at corpus scale: PubMind \citep{wang2025_pubmind} extracts variants across PubMed and PMC with Llama-3.3-70B and finds about 10\% of its mapped variants in ClinVar; on prostate-cancer literature, four LLMs extracted genes and variants with high precision and recall but occasionally introduced unsupported information \citep{Wosny2025LargeLanguage}; and AutoPM3 \citep{li2025_autopm3} combines Text2SQL over tables with retrieval over text to extract ACMG PM3 evidence, flagging supplementary tables as a hard case. When GPT-4o, Claude Opus~4 and DeepSeek-V3 extracted human disease genes from full text, disease and phenotype fields were error-prone for all three \citep{yin2025_llm_disease_genes}. A ClinGen-derived benchmark finds models that hallucinate or misinterpret results even when they classify the evidence correctly \citep{Queen2025CgbenchBenchmarking}; related work mines functional evidence \citep{saadat2026_functional_evidence} and classifies variants from literature evidence, including ClinVar's conflicting variants \citep{ma2026_aicura}. One result bears directly on our proband counting and deduplication and is unfavorable for Claude: among five models counting probands for the PS4 criterion from clinical genomics papers, Claude Sonnet~4 had the lowest exact-count concordance (73.0\%, against 91.1\% for the best model) and the lowest stability \citep{murugan2026_reasoning_llm_ps4}. Closest to our goal are literature-derived virtual cohorts built with LLMs and human curation: 1,679 Leigh syndrome spectrum cases \citep{Shen2026GenerativeAi} and about 11,000 mitochondrial disease cases \citep{Shen2026MseqdrPmdvr}.

\paragraph{Phenotypes.} HPO concept recognition spans rule- and dictionary-based tools, some with negation handling \citep{Deisseroth2019ClinphenExtracts,Liu2019Doc2hpoWeb,Groza2024FasthpocrPragmatic}, and hybrid, deep-learning and fine-tuned language-model recognizers \citep{luo2021_phenotagger,Feng2023PhenobertCombined,yang2024_phenogpt}. LLM-based recognition with GPT models was non-deterministic and needed in-context HPO labels \citep{groza2024_gpt_hpo}; GPT-4 phenotyping of OMIM clinical summaries reached concordance with manual annotators comparable to inter-rater agreement \citep{Hier2024HighthroughputPhenotyping}; retrieval-augmented RAG-HPO, benchmarked on 112 published case reports, reached a best F1 of 0.78 with Llama-3.1-70B \citep{garcia2025_raghpo}; synthetic case reports with embedding-based retrieval have been used to improve extraction \citep{Albayrak2025EnhancingHuman}; and RARE-PHENIX chains extraction, HPO standardization and prioritization on clinical notes \citep{shyr2026_rarephenix}. A study of HPO lookup through the Model Context Protocol is the source of the 97.6\% precision for Claude Sonnet~4.5 quoted in the project's application \citep{larsen2026_hpo_mcp}; that figure comes from four synthetic clinical notes in the tool-augmented condition (80.7\% without tools), with recall of 65.3\%, and Claude was not top on precision or F1, so it says little about case reports. Within the Monarch stack, SPIRES/OntoGPT \citep{caufield2024_spires} fills knowledge-base schemas zero-shot with grounding to HPO and Mondo, and GPSEA \citep{rekerle2026_gpsea} tests genotype--phenotype associations on phenopacket cohorts.

\paragraph{How GP-RareLink differs.} GP-RareLink is not a new recognizer, and we compare it with no tagger. Its contribution is the record and its evaluation: per-person, three-state records written from full-text case reports, a quote checked verbatim for every value (241,730 values with 0 span-check problems in the registered six), and ClinVar and OMIM held out of extraction so that they can score it.

\subsection{LLM-assisted systematic review and verification of claims against cited evidence}

\paragraph{Data extraction for evidence synthesis.} Most evaluations extract tens of data elements from a handful of trials or other studies, with accuracy that varies by model, prompt and field \citep{gartlehner2024_llm_extraction,Konet2024PerformanceTwo,Liu2025AidrivenEvidence,Schmidt2024ExploringUse,Mahmoudi2025CriticalAssessment,Khan2025CollaborativeLarge,kataoka2025_gpt4o_o3_extraction}; numeric fields fare worse than text \citep{kataoka2025_gpt4o_o3_extraction}, and ChatGPT-4o's reproducibility fell from 94.1\% to 77.2\% when information was not reported in the paper \citep{Jensen2025Chatgpt4oCan}, the case our ``not reported'' state is for. otto-SR automated whole reviews with 93.1\% extraction accuracy against 79.7\% for the human workflow \citep{cao2025_ottosr}. Methodological pitfalls in these evaluations have been set out \citep{Gartlehner2024FromPromise}, and the field has been reviewed systematically \citep{Shankar2026PerformanceLarge}. SLIDERS \citep{joshi2026_sliders}, from the Lam lab, adds verbatim excerpts as provenance and a code-writing reconciliation agent; the same lab's WikiChat \citep{semnani2023_wikichat} grounds outputs and verifies claims, a precedent for our gatekeeper, and STORM \citep{shao2024_storm} is its retrieval-grounded long-form synthesis system.

\paragraph{Attribution and claim verification.} AIS \citep{rashkin2023_ais} gives a formal definition of a statement supported by its cited span, the criterion our gatekeeper applies to each value and its quote. ALCE \citep{gao2023_alce} measures citation recall and precision, a human audit of generative search engines found only 51.5\% of sentences fully supported by their citations \citep{Liu2023EvaluatingVerifiability}, and FActScore \citep{min2023_factscore} checks atomic claims, the unit our audit uses (one claim per extracted value). Automatic attribution checking remains hard \citep{Honovich2022TrueReevaluating,Yue2023AutomaticEvaluation,Tang2024MinicheckEfficient}: on AttributionBench a fine-tuned GPT-3.5 reached only about 80\% macro-F1 \citep{Li2024AttributionbenchHow}. Closest to our gatekeeper, an LLM that cites evidence spans for its clinical extractions and then verifies them improved few-shot extraction \citep{gero2023_self_verification}; verification questions reduce hallucination \citep{dhuliawala2024_cove}; claim-checkers and judges for clinical text have been validated against clinician labels or a human instrument \citep{chung2025_verifact,Chung2026VerifyingFacts,croxford2025_clinical_llm_judge}; and LLM citations often fail to support the medical statements they are attached to \citep{wu2025_sourcecheckup}. Self-checking has known limits: without grounding, LLMs cannot yet self-correct their reasoning \citep{huang2024_no_self_correct}; LLM judges show self-enhancement bias \citep{zheng2023_llm_judge}, recognize and favor their own generations \citep{Panickssery2024LlmEvaluators}, and prefer text that is familiar to them \citep{Wataoka2024SelfpreferenceBias}; and LLM evaluators need human grading to stay aligned \citep{Shankar2024WhoValidates}. The application's reliability claims for Claude also need their sources' wording: its ``DRBench'' figures are hallucinated-URL rates of two search-augmented Claude models \citep{rao2026_reference_hallucination} on DeepResearch Bench \citep{du2025_deepresearchbench}, not general fabrication rates, and AA-Omniscience \citep{jackson2025_aa_omniscience} is a knowledge benchmark that rewards abstention.

\paragraph{Human audits and model comparisons.} An audit of a rare error needs care with small numerators: Wald intervals are not recommended, and score or exact intervals are used instead \citep{brown2001_binomial_ci,Wilson1927ProbableInference,Clopper1934UseConfidence}; zero and small numerators bound what a sample can certify \citep{hanley1983_zero_numerators,newman1995_small_numerators}; and agreement between raters is chance-corrected with Cohen's $\kappa$ \citep{Cohen1960CoefficientAgreement}. With 250 claims, the exact two-sided 95\% upper bound is 1.46\% when none is unsupported, 2.21\% with one and 2.86\% with two, so an audit can pass a pre-registered rule without excluding a true rate near 2\%. Closed and open-weight models have been compared across biomedical NLP benchmarks and question answering \citep{chen2025_bionlp_benchmark,stachura2025_open_weight_bioasq}; the gap closed under few-shot prompting with the same setup in report labeling \citep{dorfner2024_open_vs_commercial_radiology} and narrowed with fine-tuning in evidence summarization \citep{zhang2024_open_vs_commercial_summarization}, which is why a fair comparison fixes the scaffold. Our open-weight comparator is DeepSeek-V4.1-Flash (cf.\ the DeepSeek-V4 series; \citealp{Deepseekai2026Deepseekv4Towards}), run on the same papers with the same launch prompt and frozen files.

\paragraph{How GP-RareLink differs.} The gatekeeper judges each value strictly against its own quote, and we report its rate as a model reading, not an error rate: in the pilot, a person found 6.0\% and 6.3\% unsupported where earlier gatekeeper versions found 0.036\% and 0.34\%, and those versions flagged 2 of 6 and 1 of 4 of the claims the person rejected. The primary outcome is a pre-registered human audit of 250 claims under rule~C (at most 4 unsupported), preceded by a calibration round in which nine raters label 40 claims that carry reference labels, against a floor fixed in advance (at least 15 of 20 unsupported claims caught and 18 of 20 supported claims passed).

\subsection{Case-level rare-disease resources, imperfect reference databases, and memorization}

\paragraph{Case-level resources and benchmarks.} The phenopacket-store \citep{danis2025_phenopacket_store} is a corpus of GA4GH phenopackets hand-curated with pyphetools (6,668 phenopackets from 959 publications at v0.1.19); it is prior art, a possible human-curated reference, and a memorization risk, since phenopackets are now an evaluation currency for LLM diagnosis \citep{reese2026_llm_vs_exomiser}, and rare-disease benchmarks compile open patient datasets \citep{Chen2024RarebenchCan}. Clinical reasoning benchmarks curate items with contamination screening \citep{Shao2025MedicalagentsbenchComplex} or simulate multi-turn examinations in which LLM performance degrades \citep{Chiu2025SimulatingViva}, and leakage was the main source of high risk of bias in a meta-analysis of LLM rare-disease diagnosis \citep{nguyen2026_llm_rd_metaanalysis}. Hand curation has not yet reached most diseases: none of our twelve genes has a phenopacket-store folder.

\paragraph{Reference databases are imperfect ground truth.} ClinVar classifications disagree: literature-only entries were 8.8\% of submissions but 31.4\% of outlier classifications \citep{yang2017_clinvar_discordance}; laboratories have collaborated to resolve conflicts \citep{harrison2017_clinvar_conflicts}; 5.7\% of variants carry conflicting interpretations, mostly involving variants of uncertain significance \citep{Lazareva2024MajorCauses}; ACMG/AMP classification concordance has been measured across nine research studies \citep{Amendola2020VariantClassification}; and pathogenic ClinVar classifications have been re-examined against disease population prevalence \citep{Shah2018IdentificationMisclassified} and in the DDD study \citep{wright2021_ddd_clinvar}. The ACMG/AMP framework \citep{richards2015_acmg} and ClinGen's guidance on co-segregation and phenotype-specificity evidence \citep{biesecker2024_clingen_pp1_pp4} set how evidence from case reports is weighed when a variant is classified. Disease--phenotype annotations are incomplete as well: OMIM and Orphanet share few HPO terms for one disease \citep{connelly2026_fanconi_annotations}. When benchmark databases overlap what a method was built on, apparent accuracy is inflated \citep{Grimm2015EvaluationTools}; this circularity is why our extraction never reads ClinVar or OMIM.

\paragraph{Ascertainment bias and duplicate patients.} Literature cohorts are phenotype-first; genotype-first population analyses show lower penetrance and milder expressivity \citep{Wright2019AssessingPathogenicity,wright2024_penetrance,felker2025_allofus_genotype_first}. Spectrum estimates also depend on how absence is recorded: forms without an explicit ``absent'' option undercount features \citep{vandersluijs2019_arid1b}, which is one reason ``not reported'' must stay separate from ``excluded''. The same patients recur across sources: record linkage has deduplicated patients across independent rare-disease data sources to estimate incidence \citep{konig2019_registry_dedup}, and a few cohorts are reused across many papers \citep{shivakumar2026_cohort_reliance}, so cross-paper deduplication must be explicit.

\paragraph{Memorization.} Language models can emit verbatim training data \citep{Carlini2020ExtractingTraining,Nasr2023ScalableExtraction}, more so with scale and duplication \citep{Carlini2022QuantifyingMemorization}, and memorization has been studied in medical LLMs specifically \citep{li2026_medical_memorization}. Membership tests such as Min-K\% Prob need token log-probabilities \citep{shi2024_mink_prob}, which are available for an open-weight model but not for Claude, and apparent membership signals often reflect a shift in publication date \citep{duan2024_mia_llm}. Black-box alternatives include guided completion \citep{golchin2024_time_travel} and a contamination quiz \citep{Golchin2023DataContamination}; training cutoffs give a natural experiment \citep{roberts2024_lens_of_time}, refreshed post-cutoff items limit contamination \citep{white2025_livebench}, and the options have been surveyed \citep{xu2024_contamination_survey}.

\paragraph{How GP-RareLink differs.} ClinVar misses are triaged by people rather than all counted as extraction errors, and OMIM serves as a human-curated holdout: the registered runs recover 57.3\% (43 of 75) of the allelic variants it lists for the six diseases. A black-box memorization probe on 180 papers, stratified by publication date against Claude's training cutoff and run identically on both models, found that neither model reproduced the papers' variants or text; the after-cutoff stratum has only 10 papers, and such probes can miss memorization that changes extraction without surfacing as recall.
%%%%%%%%% RELATED WORK %%%%%%%%%

\section{Methodology}
\label{sec:method}
%%%%%%%%% METHOD %%%%%%%%%
Methodology here.
%%%%%%%%% METHOD %%%%%%%%%

\section{Experiments}
\label{sec:experiments}
%%%%%%%%% EXPERIMENTS %%%%%%%%%
Experimental setup, results, and analysis here.
%%%%%%%%% EXPERIMENTS %%%%%%%%%

\section{Conclusion}
\label{sec:conclusion}
%%%%%%%%% CONCLUSION %%%%%%%%%
Conclusion here.
%%%%%%%%% CONCLUSION %%%%%%%%%

%%%%%%%%% REFERENCES %%%%%%%%%
\bibliography{references}
\bibliographystyle{iclr2025_conference}
%%%%%%%%% REFERENCES %%%%%%%%%

%%%%%%%%% Appendix %%%%%%%%%
% If you want to add appendix, uncomment the following lines:
% \appendix
% \section{xxx}
%%%%%%%%% Appendix %%%%%%%%%

\end{document}
